\documentclass[10pt,a4paper]{amsart}
\usepackage[margin=19mm]{geometry}
\usepackage{amsmath,amssymb,mathtools}
\usepackage[scr=rsfs,cal=euler]{mathalfa}
\usepackage{xfrac}
\usepackage[unicode,hidelinks,bookmarks=false]{hyperref}
\newcommand{\ie}{i.e.\ }
\newcommand{\cf}{cf.\ }
\newcommand{\resp}{respectively\ }
\newcommand{\Subsection}{\subsection}
\providecommand{\nobreakdash}{\nobreak\hbox{-}}
\setlength{\emergencystretch}{2em}
\multlinegap0pt
\usepackage{bm}
\usepackage{bbm}
\usepackage{dsfont}
\usepackage{xspace}
\usepackage{cancel}
\usepackage{mathtools}
\usepackage{amsthm}
\makeatletter
\@ifundefined{lemma}{\newtheorem{lemma}{Lemma}}{}
\makeatother
\title[The Cauchy problem for the integrable RZQ equation]{The Cauchy problem for the integrable RZQ equation}
\author{John Holmes, Katie Massey, Ryan C. Thompson}
\date{Source: arXiv:2507.13535v2 (CC BY 4.0)}
\begin{document}
\maketitle

\noindent\emph{Source and licence.} The Cauchy problem for the integrable RZQ equation. Version arXiv:2507.13535v2 (CC BY 4.0), distributed under the Creative Commons Attribution 4.0 International licence, as recorded in the arXiv licence metadata for this exact version and shown on the abstract page https://arxiv.org/abs/2507.13535v2. Full rights notice: licenses/arxiv-2507.13535-CC-BY-4.0.txt.\par\smallskip

\noindent\textit{Editorial scope.} Counted unit: the Lemma whose source label is \texttt{lemma3} in the article (source0.tex lines 1356--1359, source label \texttt{lemma3}), delivered together with the article's own complete original proof of it (source0.tex lines 1360--1541, one \texttt{proof} environment of 182 lines). No mathematics is written, repaired or re-derived: statement and proof are the article's own text line for line. Required context retained uncounted: none. Every definition, display and label of those lines is retained unchanged. Omitted from this package and not redistributed: 1--1355, 1542--1956. Nothing inside a retained excerpt is omitted or altered: every retained body is delivered exactly as the article's lines are written, so no omission receipt of any kind accompanies this package. Citations: the retained excerpts carry no citation command at all, so no bibliography entry of the article is transcribed and no reference is redistributed. Reference adaptation: none was needed and none was made; every reference inside the delivered unit resolves inside it, so no unresolved reference or citation is produced. Printed slips kept literal: no printed word, symbol, hypothesis or number of the counted statement or of its proof was altered. Layout and class: the article's own class is \texttt{amsart} with packages including amssymb, latexsym, amsfonts, dsfont, hyperref, color, graphicx, parskip, mathrsfs, array, tikz, pgfplots; the wrapper is the lane's pinned \texttt{amsart} editorial wrapper at 10pt on A4, which restates the theorem-like environments the retained text needs and adds the article's own macro definitions that the retained text needs (none), with extra wrapper packages (bm, bbm, dsfont, xspace, cancel, mathtools). The delivered numbering is the unit's own. Assets: no retained span contains a figure environment, a graphics-inclusion command, a PDF-page inclusion, a table or a TikZ picture; no image file is copied into this package, the package declares no asset dependency and no image file is redistributed with it. Licence: version arXiv:2507.13535v2 of this article is distributed under the Creative Commons Attribution 4.0 International licence, as recorded in the arXiv licence metadata for this exact version and shown on the abstract page https://arxiv.org/abs/2507.13535v2. A full rights notice is delivered at licenses/arxiv-2507.13535-CC-BY-4.0.txt. Characters: every retained excerpt is pure ASCII, so the delivered engine, XeLaTeX with its default Latin Modern text font, reports no missing character.
\begin{lemma}\label{lemma3}
If $1 \leq \mu \leq 3, \sigma - 3 > \frac12,$ and $\mu + \sigma \geq 6,$ then 
    \[\|fg\|_{H^{\mu-3}} \leq c_{\mu,\sigma} \|f\|_{H^{\sigma-3}} \|g\|_{H^{\mu-3}}.\]
\end{lemma}

\begin{proof}
    We provide the details of Lemma \ref{lemma3} on $\mathbb{R}$, the periodic case is similar. 
%%%%%%%%%%%%%%%%%%%%%%%%%%%%%%%%%%%%%%%%%%%%%%%%%%%%%%%%%%%%%%%%%%%%%%
We first note 
\begin{equation}
    \| fg \|_{H^{\mu - 3}}^2 = \int(1 + \xi^2)^{\mu-3} \left|\widehat{fg}(\xi)\right|^2 d\xi 
\end{equation}
%%%%%%%%%%%%%%%%%%%%%%%%%%%%%%%%%%%%%%%%%%%%%%%%%%%%%%%%%%%%%%%%%%%%%%
Using the definition of convolution,
\begin{equation*}
    \| fg \|_{H^{\mu - 3}}^2 = \int(1 + \xi^2)^{\mu-3} \left|\int \hat{f}(\eta) \hat{g}(\xi - \eta)d\eta\right|^2 d\xi, 
\end{equation*}
which may be rewritten as 
\begin{equation}
    \| fg \|_{H^{\mu - 3}}^2 = 
        \int(1 + \xi^2)^{\mu-3} \left|\int (1+\eta^2)^{\frac{\sigma-3}{2}}\hat{f}(\eta)(1+\eta^2)^{\frac{3-\sigma}{2}}\hat{g}(\xi - \eta)d\eta\right|^2 d\xi. 
\end{equation}
%%%%%%%%%%%%%%%%%%%%%%%%%%%%%%%%%%%%%%%%%%%%%%%%%%%%%%%%%%%%%%%%%%%%%%
Then by Cauchy-Schwarz, 
\begin{equation}
    \| fg \|_{H^{\mu - 3}}^2 \leq 
        \|f\|_{H^{\sigma - 3}} \int(1 + \xi^2)^{\mu-3}\int(1+\eta^2)^{3-\sigma}\left| \hat{g}(\xi - \eta)\right|^2 d\eta d\xi. 
\end{equation}
%%%%%%%%%%%%%%%%%%%%%%%%%%%%%%%%%%%%%%%%%%%%%%%%%%%%%%%%%%%%%%%%%%%%%%
Implementing the change of variables $\tilde{\eta} = \xi - \eta$ and dropping the tildes, we have 
\begin{equation}
    \| fg \|_{H^{\mu - 3}}^2 \leq 
        \|f\|_{H^{\sigma - 3}} \int(1 + \xi^2)^{\mu-3}\int(1+(\xi - \eta)^2)^{3-\sigma}\left| \hat{g}(\eta)\right|^2 d\eta d\xi. 
\end{equation}
%%%%%%%%%%%%%%%%%%%%%%%%%%%%%%%%%%%%%%%%%%%%%%%%%%%%%%%%%%%%%%%%%%%%%%
A change of the order of integration then yields
\begin{equation}\label{above}
    \| fg \|_{H^{\mu - 3}}^2 \leq 
        \|f\|_{H^{\sigma - 3}} \int \left| \hat{g}(\eta)\right|^2 \int(1 + \xi^2)^{\mu-3}(1+(\xi - \eta)^2)^{3-\sigma} d\xi d\eta. 
\end{equation}
%%%%%%%%%%%%%%%%%%%%%%%%%%%%%%%%%%%%%%%%%%%%%%%%%%%%%%%%%%%%%%%%%%%%%%
For the integral in \eqref{above} to be bounded by $c_{\mu,\sigma}\|g\|_{\mu-3},$ it remains to show
\begin{equation}\label{wwts}
    I(\eta) := \int(1 + \xi^2)^{\mu-3}(1+(\xi - \eta)^2)^{3-\sigma}                 d\xi 
            \leq c_{\mu,\sigma}(1+\eta^2)^{\mu-3}.
\end{equation}
%%%%%%%%%%%%%%%%%%%%%%%%%%%%%%%%%%%%%%%%%%%%%%%%%%%%%%%%%%%%%%%%%%%%%%
After making the change of variables $\tilde{\xi} = \xi - \eta$, dropping the tilde, and rearranging, we arrive at 
\begin{equation}
    I(\eta) = \int \frac{d \xi}{(1+(\xi - \eta)^2)^{3-\mu}(1 + \xi^2)^{\sigma-3}}.
\end{equation}
%%%%%%%%%%%%%%%%%%%%%%%%%%%%%%%%%%%%%%%%%%%%%%%%%%%%%%%%%%%%%%%%%%%%%%
Here, we note that since $3 - \mu > 0$ and $\sigma - 3 > \frac12,$ the integral is finite for all $\eta.$ 
%%%%%%%%%%%%%%%%%%%%%%%%%%%%%%%%%%%%%%%%%%%%%%%%%%%%%%%%%%%%%%%%%%%%%%
Additionally, we may consider only the case in which $\eta \geq 0$ since when $\eta < 0,$ the change of variables $\tilde{\xi} = -\xi$ returns the problem to the $\eta \geq 0$ case.  
Lastly, for $\eta >0 $ and $\xi \leq 0$, we have $\xi^2 \geq 2\xi\eta$ which gives $(1+(\xi - \eta)^2)^{3-\mu}\geq(1 + \eta^2)^{3-\mu}$. 
%%%%%%%%%%%%%%%%%%%%%%%%%%%%%%%%%%%%%%%%%%%%%%%%%%%%%%%%%%%%%%%%%%%%%%
Thus \eqref{wwts} holds on $(-\infty, 0]$ and we are left to show 
\begin{equation}
    J(\eta) := \int_0^{\infty}\frac{d \xi}{(1+(\xi - \eta)^2)^{3-\mu}(1 + \xi^2)^{3-\sigma}} \leq \frac{c_{\mu,\sigma}}{(1 + \eta^2)^{3-\mu}}, \ \ \eta>0.
\end{equation}
%%%%%%%%%%%%%%%%%%%%%%%%%%%%%%%%%%%%%%%%%%%%%%%%%%%%%%%%%%%%%%%%%%%%%%
We will estimate by breaking the domain of integration into the subintervals 
    $[0, \frac\eta2], [\frac\eta2,\eta], [\eta,\frac{3\eta}{2}],$ $ [\frac{3\eta}{2},\infty)$. 
%%%%%%%%%%%%%%%%%%%%%%%%%%%%%%%%%%%%%%%%%%%%%%%%%%%%%%%%%%%%%%%%%%%%%%
For $\xi \in [0,\frac\eta2 ],$ we have 
    $1 + (\xi - \eta)^2 \geq 1 + \frac{\eta^2}{4}$ and thus 
\begin{equation*}
    \int_0^{\frac{\eta}{2}}\frac{d \xi}{(1+(\xi - \eta)^2)^{3-\mu}(1 + \xi^2)^{\sigma-3}} 
        \leq \frac{1}{(1+\eta^2/4)^{3-\mu}}\int_0^{\frac{\eta}{2}}\frac{d \xi}{(1 + \xi^2)^{\sigma-3}}.
\end{equation*}
%%%%%%%%%%%%%%%%%%%%%%%%%%%%%%%%%%%%%%%%%%%%%%%%%%%%%%%%%%%%%%%%%%%%%%
Since $\mu - 3 \leq 0$, we see 
    $(1 + \frac{\eta^2}{4})^{\mu-3} = \left[\frac14( 4 + \eta^2)\right]^{\mu-3} \leq 4^{3-\mu}(1 + \eta^2)^{\mu-3}$.
%%%%%%%%%%%%%%%%%%%%%%%%%%%%%%%%%%%%%%%%%%%%%%%%%%%%%%%%%%%%%%%%%%%%%%
Therefore 
\begin{equation}
    \int_0^{\frac{\eta}{2}}\frac{d \xi}{(1+(\xi - \eta)^2)^{3-\mu}(1 + \xi^2)^{\sigma-3}} 
        \leq \frac{4^{3-\mu}c_s}{(1+\eta^2)^{3-\mu}},
\end{equation}
where 
    \[c_s = \int_0^{\infty}\frac{d \xi}{(1 + \xi^2)^{\sigma-3}}.\]
%%%%%%%%%%%%%%%%%%%%%%%%%%%%%%%%%%%%%%%%%%%%%%%%%%%%%%%%%%%%%%%%%%%%%%
The interval $\xi \in \left[\frac{3\eta}{2}, \infty
\right)$ can be estimated in the same manner since the inequality $1 + (\xi - \eta)^2 \geq 1 + \frac{\eta^2}{4}$ still holds. 
%%%%%%%%%%%%%%%%%%%%%%%%%%%%%%%%%%%%%%%%%%%%%%%%%%%%%%%%%%%%%%%%%%%%%%
As such, 
\begin{equation}
    \int_{\frac{3\eta}{2}}^{\infty}\frac{d \xi}{(1+(\xi - \eta)^2)^{3-\mu}(1 + \xi^2)^{\sigma-3}} 
        \leq \frac{4^{3-\mu}c_s}{(1+\eta^2)^{3-\mu}},
\end{equation}
%%%%%%%%%%%%%%%%%%%%%%%%%%%%%%%%%%%%%%%%%%%%%%%%%%%%%%%%%%%%%%%%%%%%%%
For $\xi \in \left[\frac\eta2,\eta\right],$ we have $(1 + \xi^2)^{\sigma - 3} \geq (1 + \frac{\eta^2}{4})^{\sigma-3}$. 
%%%%%%%%%%%%%%%%%%%%%%%%%%%%%%%%%%%%%%%%%%%%%%%%%%%%%%%%%%%%%%%%%%%%%%
We additionally note that for $a \geq 0, 1 + a^2 \geq \frac12(1+a)^2.$
%%%%%%%%%%%%%%%%%%%%%%%%%%%%%%%%%%%%%%%%%%%%%%%%%%%%%%%%%%%%%%%%%%%%%%
Letting $a = \eta - \xi$ and raising to the $3 - \mu$ power yields
    $(1 + (\xi - \eta)^2)^{3-\mu} \geq 2^{\mu - 3}(1 + \eta - \xi)^{2}$. 
%%%%%%%%%%%%%%%%%%%%%%%%%%%%%%%%%%%%%%%%%%%%%%%%%%%%%%%%%%%%%%%%%%%%%%
Then 
\begin{equation}\label{J3}
    \int_{\frac{\eta}{2}}^{\eta}\frac{d \xi}{(1+(\xi - \eta)^2)^{3-\mu}(1 + \xi^2)^{\sigma-3}} 
        \leq \frac{2^{3-\mu}4^{\sigma-3}}{(1+\eta^2)^{\sigma-3}}\int_{\frac{\eta}{2}}^{\eta}\frac{d\xi}{(1+\eta-\xi)^{2(3-\mu)}}.
\end{equation}
%%%%%%%%%%%%%%%%%%%%%%%%%%%%%%%%%%%%%%%%%%%%%%%%%%%%%%%%%%%%%%%%%%%%%%
It remains to show 
\begin{equation*}
    K(\eta) := 
         \int_{\frac{\eta}{2}}^{\eta}\frac{d\xi}{(1+\eta-\xi)^{2(3-\mu)}} < \infty.
\end{equation*}
%%%%%%%%%%%%%%%%%%%%%%%%%%%%%%%%%%%%%%%%%%%%%%%%%%%%%%%%%%%%%%%%%%%%%%
We consider three cases: $\mu \in \left[1,\frac52\right), \mu = \frac52,$ and $\mu \in \left(\frac52, 3\right]$. 
%%%%%%%%%%%%%%%%%%%%%%%%%%%%%%%%%%%%%%%%%%%%%%%%%%%%%%%%%%%%%%%%%%%%%%
For $\mu \in \left[1,\frac52\right)$, we have $2\mu-5 < 0$ and so evaluating the integral gives 
\begin{equation}\label{K1}
    K(\eta) = 
         \frac{1}{5-2\mu}\left[1-(1+\eta/2)^{2\mu-5}\right] 
         \leq \frac{1}{5-2\mu}.
\end{equation}
%%%%%%%%%%%%%%%%%%%%%%%%%%%%%%%%%%%%%%%%%%%%%%%%%%%%%%%%%%%%%%%%%%%%%%
Combining \eqref{J3} and \eqref{K1} and using the fact that $3-\sigma \leq \mu - 3$, we have
\begin{equation}
    \int_{\frac{\eta}{2}}^{\eta}\frac{d \xi}{(1+(\xi - \eta)^2)^{3-\mu}(1 + \xi^2)^{\sigma-3}} 
        \leq \frac{2^{3-\mu}4^{\sigma-3}}{(5-2\mu)(1+\eta^2)^{\sigma-3}}
        \leq \frac{2^{3-\mu}4^{\sigma-3}}{(5-2\mu)(1+\eta^2)^{3-\mu}}.
\end{equation}
%%%%%%%%%%%%%%%%%%%%%%%%%%%%%%%%%%%%%%%%%%%%%%%%%%%%%%%%%%%%%%%%%%%%%%
For $\mu \in \left(\frac52,3\right]$, we have $2\mu-5 > 0$ and thus
\begin{equation*}
    K(\eta) = 
         \frac{1}{2\mu-5}\left[(1+\eta/2)^{2\mu-5})-1\right] 
         \leq \frac{1}{2\mu-5}\left(1+\eta/2\right)^{2\mu-5}.
\end{equation*}
%%%%%%%%%%%%%%%%%%%%%%%%%%%%%%%%%%%%%%%%%%%%%%%%%%%%%%%%%%%%%%%%%%%%%%
Noting that for $a \geq 0, \frac12(1+a)^2 \leq 4 + a^2$ and substituting in $a = \frac\eta2,$ we find 
    $(1 + \frac\eta2)^{2\mu-5} \leq 2^{\mu-5/2}(1 + \eta^2)^{\mu-5/2}.$ 
%%%%%%%%%%%%%%%%%%%%%%%%%%%%%%%%%%%%%%%%%%%%%%%%%%%%%%%%%%%%%%%%%%%%%%
The above estimate then becomes
\begin{equation}\label{K2}
    K(\eta) \leq  
         \frac{2^{\mu-5/2}(1 + \eta^2)^{\mu-5/2}}{2\mu-5}.
\end{equation}
%%%%%%%%%%%%%%%%%%%%%%%%%%%%%%%%%%%%%%%%%%%%%%%%%%%%%%%%%%%%%%%%%%%%%%
Combining \eqref{J3} and \eqref{K2} and using the fact that $(\sigma-3) - (\mu - 5/2) \geq 3 - \mu$, we have
\begin{equation}
    \int_{\frac{\eta}{2}}^{\eta}\frac{d \xi}{(1+(\xi - \eta)^2)^{3-\mu}(1 + \xi^2)^{\sigma-3}} 
      %  \leq \frac{2^{3-\mu}4^{\sigma-3}2^{\mu-5/2}(1 + \eta^2)^{\mu-5/2}}{(2\mu-5)(1+\eta^2)^{\sigma-3}}
        \leq \frac{2^{1/2}4^{\sigma-3}}{(2\mu-5)(1+\eta^2)^{3-\mu}}.
\end{equation}
%%%%%%%%%%%%%%%%%%%%%%%%%%%%%%%%%%%%%%%%%%%%%%%%%%%%%%%%%%%%%%%%%%%%%%
Finally, if $\mu = \frac52,$ 
\begin{equation}\label{K3}
    K(\eta) = 
        \int_{\frac{\eta}{2}}^{\eta}\frac{d\xi}{(1+\eta-\xi)}
        = \ln\left(1 + \frac\eta2\right).
\end{equation}
%%%%%%%%%%%%%%%%%%%%%%%%%%%%%%%%%%%%%%%%%%%%%%%%%%%%%%%%%%%%%%%%%%%%%%
Since $\ln\left(1 + \frac\eta2\right) \leq c_{\sigma}'\left(1 + \frac\eta2\right)^{\sigma-7/2}$, we have by combining \eqref{J3} and \eqref{K3}, 
\begin{equation}
    \int_{\frac{\eta}{2}}^{\eta}\frac{d \xi}{(1+(\xi - \eta)^2)^{3-\mu}(1 + \xi^2)^{\sigma-3}} 
       % \leq \frac{2^{3-\mu}4^{\sigma-3}c_{\sigma}(1 + \eta^2)^{\sigma-7/2}}{(1+\eta^2)^{\sigma-3}}
        \leq \frac{2^{3-\mu}4^{\sigma-3}c_{\sigma}}{(1+\eta^2)^{3-\mu}}.
\end{equation}
Therefore for $\mu \in \left[1,3\right],$
\begin{equation}
    \int_{\frac{\eta}{2}}^{\eta}\frac{d \xi}{(1+(\xi - \eta)^2)^{3-\mu}(1 + \xi^2)^{\sigma-3}} 
        \leq \frac{c_{\mu,\sigma}}{(1+\eta^2)^{3-\mu}}.
\end{equation}
%%%%%%%%%%%%%%%%%%%%%%%%%%%%%%%%%%%%%%%%%%%%%%%%%%%%%%%%%%%%%%%%%%%%%%
For $\xi \in \left[\eta,\frac{3\eta}{2}\right]$, we note that $(1+\xi^2)^{\sigma - 3} \geq (1 + \eta^2)^{\sigma - 3}$. 
%%%%%%%%%%%%%%%%%%%%%%%%%%%%%%%%%%%%%%%%%%%%%%%%%%%%%%%%%%%%%%%%%%%%%%
Using the same identity as in \eqref{J3} now with $a = \xi - \eta$, we have $(1 + (\xi - \eta)^2) \geq 2^{\mu-3}(1 + \xi - \eta)^2$. 
%%%%%%%%%%%%%%%%%%%%%%%%%%%%%%%%%%%%%%%%%%%%%%%%%%%%%%%%%%%%%%%%%%%%%%
Thus
\begin{equation}\label{J4}
    \int_{\eta}^{\frac{3\eta}{2}}\frac{d \xi}{(1+(\xi - \eta)^2)^{3-\mu}(1 + \xi^2)^{\sigma-3}} 
        \leq \frac{2^{3-\mu}}{(1+\eta^2)^{\sigma-3}}\int^{\frac{3\eta}{2}}_{\eta}\frac{d\xi}{(1+\xi-\eta)^{2(3-\mu)}}.
\end{equation}
%%%%%%%%%%%%%%%%%%%%%%%%%%%%%%%%%%%%%%%%%%%%%%%%%%%%%%%%%%%%%%%%%%%%%%
Remaining calculations follow similarly to the above case, and we again obtain the bound 
\begin{equation}
    \int_{\eta}^{\frac{3\eta}{2}}\frac{d \xi}{(1+(\xi - \eta)^2)^{3-\mu}(1 + \xi^2)^{\sigma-3}} 
        \leq \frac{c_{\mu,\sigma}}{(1+\eta^2)^{3-\mu}}.
\end{equation}
%%%%%%%%%%%%%%%%%%%%%%%%%%%%%%%%%%%%%%%%%%%%%%%%%%%%%%%%%%%%%%%%%%%%%%
Thus the proof of Lemma \ref{lemma3} on $\mathbb{R}$ is complete.
\end{proof}

\end{document}
